\textbf{The system has not been evaluated.} No human-subjects study was conducted, and no
end-to-end integration testing of the complete pipeline has been carried out. Detection accuracy,
sensor calibration, interruption load, latency, and the effectiveness of the intervention layer are
all unmeasured. The thresholds, hold times and severity rules in Table~\ref{tab:parameters} are
design values chosen from published ergonomic guidance, not values validated against observation.
This is the principal limitation, it bounds all the others, and it means the contribution of this
paper is an architecture and its implementation rather than an empirical result.

The remaining limitations fall into three groups.

\subsection{Measurement}

\begin{itemize}
    \item \textbf{Head angles are relative, not anatomical.} Orientation is measured against a
    user-set zero captured at calibration, so values express deviation from an individual's chosen
    neutral rather than a cervical angle. The relationship between visible head posture and symptoms
    is in any case not deterministic: Sarig Bahat et al. found no clinically applicable association
    between forward-head posture and pain intensity~\cite{ref22}. Postural output should be read as
    an exposure indicator, never as evidence of a disorder.
    \item \textbf{A front-facing monocular camera cannot recover sagittal depth.} Forward-head
    displacement, trunk flexion and slouching are excluded for this reason
    (Table~\ref{tab:posture_variables_detail}), so the postural picture is partial by construction.
    \item \textbf{The two distance channels reference different points.} Their agreement bounds
    precision; it does not establish accuracy for either.
    \item \textbf{Illuminance is derived, not measured.} The ADC-to-lux conversion is a modelled
    curve over a photoresistor divider whose coefficients are nominal rather than fitted against a
    photometric reference. Above the divider's saturation point the reading is censored --- reported
    as a lower bound rather than a value. With the fixed resistor used here, the trustworthy ceiling
    falls close to the lower edge of the target illuminance band, so ordinary desk lighting can push
    the converter beyond its linear region.
    \item \textbf{A single illuminance sensor cannot characterise lighting quality.} Glare,
    reflections, light direction, screen luminance and colour temperature all matter for visual
    comfort~\cite{ref7} and none are observed. Ambient light is a contextual variable only.
    \item \textbf{Blink rate is not an independent diagnostic.} It varies with task demand and
    cognitive load~\cite{ref11}, and is used here as a supporting indicator.
\end{itemize}

\subsection{Scope of Operation}

\begin{itemize}
    \item \textbf{Exposure is counted only while a face is detected}, so poor lighting or occlusion
    cause under-counting in exactly the conditions the system exists to flag.
    \item \textbf{Exposure is continuous, not cumulative.} The measured quantity is uninterrupted
    desk time within a session and does not estimate daily sedentary load, so it cannot be compared
    against the daily-duration odds ratios cited in Section~\ref{sec:litreview}.
    \item \textbf{Thresholds and hold times are derived from published guidance and from
    calibration, not from validation against ergonomic assessment by an expert.} They are plausible
    defaults, not established cut-points.
    \item \textbf{The recommendation layer degrades to fixed phrases} when the language model API is
    unreachable, and generated and templated output are not equivalent.
    \item \textbf{Camera placement, occlusion and lighting influence computer-vision
    performance}~\cite{ref16}, and the streamed resolution is well below the sensor's capability.
\end{itemize}

\subsection{Deployment}

\begin{itemize}
    \item \textbf{Image data leaves the device.} Vision processing is performed off the
    microcontroller, so any real deployment must address retention, access and transmission
    security for facial video --- a requirement an on-device inference design would have avoided.
    \item \textbf{No long-term adherence data.} Whether the governor's design actually prevents the
    alert fatigue it is intended to prevent has not been observed over any meaningful period.
    \item \textbf{The system provides ergonomic risk screening and behavioural support.} It does not
    diagnose musculoskeletal disorders, Computer Vision Syndrome, or any other condition, and its
    recommendations have not been reviewed by ergonomics or healthcare professionals.
\end{itemize}
